\section{Scope, Method, and Qualifications}

\subsection{Reader, question, and method}

This is a discursive theological article for a serious general reader, able to follow careful distinctions but not presumed to know scholastic terminology. It is Part III, the conclusion, of a three-part series, \emph{The Creature Before God}; each part is a complete essay. Part I (\emph{Ontological Vertigo}) established received being, present conservation, the one-sided real relation, and the fork of recoil or consent; Part II (\emph{The Due Return}) established the unpayable debt, the virtue of religion, and the open account, ending at an arrow from gratitude toward charity. This essay follows that arrow. Its governing question is what the creature is finally for: what stands at the end of every why once justice has named the necessary rendering. Its claim: the will of the rational creature is calibrated to the universal good, which no created good is; the only end that would close its desire is the vision of the divine essence, which no created power can attain; this incapacity is the creature's rank, not its defect; that the end is nevertheless offered is a datum of revelation, not of the ascent---God has ordered man to a supernatural end, and the vision is dogmatically guaranteed, no creature mediating; the love by which the creature moves toward that end is charity, friendship with God, founded on God's communication of His own happiness and infused, not acquired; hope is the virtue of the interval, leaning on the divine help for nothing less than God Himself; and the consummation---rest, vision, love, praise, in an account open forever---perfects rather than repeals the creature's dependence, freedom, and due return. The trilogy thus closes where it began: the falling that was being carried arrives at the face of the carrier.

The argument proceeds in epistemic order: the last-end structure of desire (ST I-II, q.~1), stated without theological premises beyond the participation framework adopted in Part I; the audit of created goods and the universal-good argument (I-II, q.~2); the identification of the satisfying end as vision (I-II, q.~3, a.~8; I, q.~12, a.~1) with the incapacity of natural powers (I, q.~12, a.~4; I-II, q.~5, a.~5) and the dignity of the capacity (I-II, q.~5, a.~5, ad~2); the seam to revelation expressly marked where the ascent runs out; the offered end (\emph{Dei Filius} ch.~2; 1 Cor 13:12; 1 Jn 3:1--2), the light of glory (I, q.~12, aa.~2, 4--5), and the dogmatic definition (\emph{Benedictus Deus}; CCC 1023); charity as friendship (John 15:15; II-II, q.~23, a.~1; q.~24, a.~2) with its boundaries; hope (II-II, q.~17, aa.~1--2) between presumption and despair; and the consummation (\emph{De civitate Dei} XXII.30; CCC 1023--1029, 1718--1729) with the trilogy's deliberate close. Revealed claims are labeled as received where they appear and are nowhere presented as conclusions of the philosophical thread. English is the working language; no liturgical rite, civil or canonical jurisdiction, or mutable discipline governs any conclusion. All online witnesses were checked on 2026-07-25; the City of God quotations were checked in the source library's tracked 1871 transcription at its registered passage.

\begin{threestable}{Source class}{Function in the article}{Governing boundary}
Scripture & John 15:9--16 supplies the friendship datum; 1 Cor 13:8--13 the mirror, the remaining three, and charity's permanence; 1 Jn 3:1--2 filiation and vision; Rom 5:5 infusion; Ps 26:8 the sought face; 1 Cor 2:9 within \emph{Dei Filius}. & Quoted from the public-domain Douay--Rheims (Challoner) at exact loci; readings are textual and modest; no independent exegetical apparatus is claimed.\\
Thomistic framework & ST I-II, qq.~1--5; I, q.~12; II-II, qq.~17, 23--24 govern the analysis of the last end, the vision, hope, and charity; II-II, q.~106 and Part I--II loci are inherited. & Adopted as the article's framework where Catholic theology permits schools; not asserted as dogma; each locus used at its exact argumentative level.\\
Patristic witness & Augustine, \emph{Confessions} I.1.1 supplies the first-person register of the restless heart; \emph{De civitate Dei} XXII.30 the consummation: God as self-promised reward, partakers not gods, perfected freedom, the sabbath. & Cited as witness and record, not demonstration; the Dods 1871 edition witness is identified and its wording verified in the tracked source-library transcription.\\
Authoritative Catholic teaching & \emph{Dei Filius} ch.~2 (supernatural end); \emph{Benedictus Deus} (the beatific vision defined); CCC 1023--1029, 1718--1729 (heaven; the vocation to beatitude). & Cited as the Church's teaching at its stated authority, explicitly beyond what the philosophy proved; dogmatic force confined to what the texts state.\\
Project synthesis & The why-chain opening and its three endings; the corridor/room, compass, ceiling, open-mouth, scaffolding, and letters-from-the-road images; the mapping of presumption and despair onto Parts I--II's coinages; the reading of the open account as friendship's circulation; all final formulations. & Original editorial synthesis argued from the checked sources; attributed to no source and creating no doctrine.\\
\end{threestable}

\subsection{Included and excluded scope}

Included are: the necessity of a last end in the order of intention (I-II.1.4, 1.6) and the formal agreement of all in desiring it (1.7); the audit of created candidate goods (2.1--7) and the universal-good conclusion (2.8); beatitude as vision of the divine essence (3.8) with the natural-desire argument (I.12.1) at its stated level; the incapacity of natural powers (I.12.4; I-II.5.5) and the nobility of assisted capacity (5.5 ad~2); the supernatural end as dogma (\emph{Dei Filius} ch.~2); the promise texts (1 Cor 13:12; 1 Jn 3:1--2; Ps 26:8); grace-union and the light of glory (I.12.2, 4--5); the definition of \emph{Benedictus Deus} with CCC 1023; charity as friendship founded on communicated beatitude (II-II.23.1 with ad~1), its infusion (24.2), its primacy and form-giving role (23.6, 23.8, cited summarily); the landing of Part II's gratitude arrow (106.6 ad~2); hope's object and leaning (17.1--2) with presumption and despair named structurally; the permanence of charity (1 Cor 13:8, 13); and the consummation (De civ. Dei XXII.30; CCC 1020s and 1720s) with the trilogy's close.

Excluded are: treatise-scale accounts of grace, justification, merit, predestination, assurance, and perseverance; the mechanics of infused virtue and the gifts of the Holy Spirit; Christology, soteriology, and the Passion (named as presupposed by the tradition's texts, nowhere developed); ecclesiology and the sacraments, including Eucharistic theology and the sacramental economy of the meantime; purgatory, judgment, hell, and the resurrection of the body beyond the definition's own boundary; the timing and states of eschatology beyond \emph{Benedictus Deus}'s stated scope; mystical theology, contemplative prayer, private revelation, and any phenomenology of religious experience; the scholastic disputes over the natural desire for the supernatural (named, not adjudicated); pastoral counsel on despair, scrupulosity, or assurance; and judgments about any person's state, culpability, or salvation.

\subsection{Material qualifications}

\begin{enumerate}[leftmargin=1.45em,itemsep=0.25em]
\item \emph{Ontological vertigo}, \emph{recoil}, \emph{consent} (Part I), \emph{the due return}, \emph{the open account}, and \emph{spiritual bookkeeping} (Part II) are this series' coinages, used in their defined senses; \emph{beatitude}/\emph{happiness} translates \latin{beatitudo}, the complete good, not a mood.
\item The last-end and universal-good arguments presuppose the participation metaphysics adopted and bounded in Part I; a reader who rejects that framework is not offered independent grounds here.
\item ST I-II.3.8 and I.12.1 stand in a work of Christian theology; the essay uses them for the conditional (nothing less than vision would satisfy) and flags that ``the blessed see the essence of God'' is Aquinas's theological conclusion. How much the natural-desire argument proves is contested among Thomists; this essay does not adjudicate the dispute, and nowhere treats the desire as an exigency binding God.
\item The gratuity of the supernatural end is dogmatically protected (\emph{Dei Filius} ch.~2: \latin{ex infinita bonitate sua}); the essay's ``question addressed to God'' image asserts an openness of the creature, not a claim on the answer.
\item \emph{Benedictus Deus} defines the state of purified souls before the resurrection; the essay confines its use to the definition's content (intuitive, facial vision, no creature mediating) and takes no position on disputed eschatological questions outside it.
\item The light of glory is a created disposition of the seer, not a created object seen; ``made deiform'' (I.12.5) and ``partakers of the divine nature'' (2 Pet 1:4 in CCC 1721) assert participation, never identity of essence, absorption, or loss of creaturehood.
\item Charity as friendship (II-II.23.1) rests on God's communication of His own beatitude; the essay derives no mutuality from creaturely resources and retains Part I's one-sided real relation intact. Friendship does not abolish \latin{latria}, obligation, or the Creator--creature asymmetry.
\item The infusion of charity and its identification with participation in the Holy Spirit's love are received doctrine (II-II.24.2 with Rom 5:5); the Trinitarian identification is revelation, marked as beyond the ascent in Part I and not argued here.
\item Presumption and despair are treated structurally as the vices opposed to hope; the mapping onto ``spiritual bookkeeping'' and ``recoil'' is project synthesis, and no claim is made about the culpability or interior state of any person; pastoral treatment of despair is outside scope.
\item The reading of 1 Cor 13 (faith and hope as proper to the interval, charity permanent) follows the common doctrinal reading; the scholastic details of the cessation of faith and hope in glory are not developed.
\item Augustine's XXII.30 is quoted from the Dods 1871 translation; the Catechism's own rendering of the same sentences (CCC 1720) differs verbally and both witnesses are identified exactly; neither is altered or conflated.
\item The perfected freedom of the blessed (``not able to sin'') is Augustine's teaching in XXII.30, received by the broad tradition; the essay uses it against the rivalry picture diagnosed in Part I and makes no claim about the mechanics of impeccability.
\item The essay's closing sentences are deliberate synthesis closing the series; ``pardon\ldots{} another story, told elsewhere'' marks redemption as presupposed and undeveloped, not as unnecessary.
\item No pastoral, canonical, or liturgical advice is given; nothing here diagnoses any person's hope, charity, or destiny.
\end{enumerate}

\subsection{Notes, translations, rights, currentness, and review}

The numbered Notes carry exact citations, edition and translation identifications, verification dates, and qualifications that would interrupt the argument; no indispensable premise lives only in a note. Scripture is quoted from the public-domain Douay--Rheims translation (Challoner revision) at drbo.org. Aquinas is quoted from the public-domain English Dominican translation with Latin checked at Corpus Thomisticum. Augustine's \emph{Confessions} is quoted from the public-domain Pilkington translation (NPNF~1/1) with the Latin phrase checked at a public Latin text; \emph{De civitate Dei} from the public-domain Dods translation (1871), checked in the source library's tracked transcription of the 1871 printing at the registered passage for XXII.30. \emph{Dei Filius} is quoted from the Holy See's Latin web text with labeled working glosses; \emph{Benedictus Deus} from the received Denzinger Latin (DH 1000--1002, checked at a public bilingual transcription) with its central definition in the official English of CCC 1023 and remaining renderings as labeled glosses; the Catechism is quoted briefly from the official Vatican English web text with attribution. Official texts and third-party translations remain outside the project's CC~BY~4.0 grant; project-created prose, images, and organization are project content.

All online witnesses were checked on 2026-07-25; no mutable law or discipline is relied on. This revision received internal argumentative, source-consistency, quotation, rights, and production review by the authoring agent. Independent scriptural, patristic, Thomistic, dogmatic, theological, literary, and ecclesiastical review is outstanding; no imprimatur, nihil obstat, or ecclesiastical approval is claimed, and internal checking is not independent review.
